\documentclass[11pt]{article}

\usepackage[T1]{fontenc}
\usepackage{amsmath,amssymb,amsthm}
\usepackage[margin=1in]{geometry}
\usepackage[hidelinks]{hyperref}

\newtheorem{theorem}{Theorem}
\newtheorem{lemma}[theorem]{Lemma}
\newtheorem{corollary}[theorem]{Corollary}
\theoremstyle{remark}
\newtheorem{remark}[theorem]{Remark}

\newcommand{\G}{\mathcal G}
\newcommand{\Z}{\mathbb Z}

\title{A Sharp Width-Ten Gap for Defective Valid-Extension Sets\\
of Finite Gilbreath Sequences}
\author{Xiaojian Zeng\\
PBC School of Finance, Tsinghua University, 100190\\
\texttt{xiaojianzeng@foxmail.com}}
\date{}

\hypersetup{
  pdftitle={A Sharp Width-Ten Gap for Defective Valid-Extension Sets of Finite Gilbreath Sequences},
  pdfauthor={Xiaojian Zeng}
}

\begin{document}

\maketitle

\begin{abstract}
For a normalized strictly increasing Gilbreath sequence $S$, let $K_S$ be
the full two-sided set of integers that can be appended while preserving the
Gilbreath property, and let $C_S$ be its natural parity interval.  We prove
the all-length implication
\[
  K_S\ne C_S \quad\Longrightarrow\quad |K_S|\ge 10.
\]
The bound is sharp at $S=(2,3,5,9,15)$.  Consequently, extension width at
most nine is equivalent to interval-completeness together with an
anti-diagonal sum at most seven.  This gives four possible ordered
anti-diagonal types.  It does not classify the top rows realizing those
types, so the original all-length stability-classification problem remains
unresolved.
\end{abstract}

\section{Context and statement}

For an integer sequence $S=(s_1,\ldots,s_n)$, define its absolute-difference
triangle by
\[
 s_a^0=s_a,\qquad
 s_a^b=\lvert s_{a+1}^{b-1}-s_a^{b-1}\rvert
 \quad(1\le b\le n-1,\ 1\le a\le n-b).
\]
The sequence is \emph{Gilbreath} when $s_1^b=1$ for every positive row.
Write
\[
 \G_n=\{(s_1,\ldots,s_n):2=s_1<s_2=3<s_3<\cdots<s_n,
               \ s_1^b=1\ (1\le b<n)\}.
\]
For $S\in\G_n$, its full valid-extension set is
\[
 K_S=\{k\in\Z:(s_1,\ldots,s_n,k)\text{ is Gilbreath}\}.
\]
There is no monotonicity requirement on the appended integer $k$.

Muney \cite{Muney2026} developed an exact theory of these sets.  For the
ordered right anti-diagonal
\[
 e_i=s_{n-i}^i\quad(1\le i<n),\qquad
 A(S)=\sum_{i=1}^{n-1}e_i,
\]
the natural candidate set is
\[
 C_S=\{k\in\Z:|k-s_n|\le A(S)+1,
                     \ k\equiv s_n\pmod 2\}.
\]
The closest results needed here are Muney's reverse-tree characterization
and ordered interval-completeness criterion \cite[Proposition~18 and
Theorem~20]{Muney2026}.  The latter is an order-sensitive counterpart of
the classical subset-sum completeness criterion of Brown
\cite{Brown1961}.  Muney also found the first failure of $K_S=C_S$, proved
that the minimum extension width is five, and asked for an all-length
classification of the members of $\G_n$ having width at most nine
\cite[Theorems~24--25 and Question~16(5)]{Muney2026}.

Our contribution is a sharp gap theorem for the defective case
$K_S\ne C_S$, followed by a finite reduction in the language of ordered
anti-diagonals.  Muney's paper does not state this gap or the reduction.
We make no broader priority claim.  In particular, the reduction is not the
top-row classification requested in the source problem.

\section{Reverse sets and realizability constraints}

For $e\ge0$ and a finite $T\subseteq\Z_{\ge0}$, put
\[
 P_e(T)=\{e+t:t\in T\}\cup\{e-t:t\in T,\ e\ge t\}.
\]
The condition $e\ge t$ in the second branch is essential.  Starting with
$T_{n-1}=\{1\}$, define
\[
 T_{i-1}=P_{e_i}(T_i)\qquad(i=n-1,n-2,\ldots,1).
\]
Muney's reverse-tree theorem gives
\begin{equation}\label{eq:reverse}
 D_S:=\{|k-s_n|:k\in K_S\}=T_0,
 \qquad
 K_S=\{s_n+d:d\in D_S\}\cup\{s_n-d:d\in D_S\}.
\end{equation}
The interval-completeness criterion is
\begin{equation}\label{eq:criterion}
 K_S=C_S
 \quad\Longleftrightarrow\quad
 e_i\le 1+\sum_{j>i}e_j
 \quad(1\le i\le n-2).
\end{equation}

The proof of the gap uses two constraints that distinguish realizable
anti-diagonals from arbitrary folding tuples.

\begin{lemma}[Parity]\label{lem:parity}
If $S\in\G_n$, then every $s_j$ with $j\ge2$ is odd.  Consequently
\[
 e_1,\ldots,e_{n-2}\text{ are even},\qquad e_{n-1}=1.
\]
\end{lemma}

\begin{proof}
The assertion starts with $s_2=3$.  For $j\ge3$, reduce the bottom entry of
the triangle of the prefix $(s_1,\ldots,s_j)$ modulo two.  Since absolute
subtraction is addition modulo two,
\[
 1\equiv \sum_{a=1}^{j}\binom{j-1}{a-1}s_a\pmod2.
\]
Here $s_1=2$ contributes zero.  If $s_2,\ldots,s_{j-1}$ are odd, their total
contribution is
\[
 \sum_{a=2}^{j-1}\binom{j-1}{a-1}=2^{j-1}-2\equiv0\pmod2.
\]
Thus $s_j$ is odd, completing the induction.

The first positive row therefore begins with $1$ and has only even entries
thereafter.  The same description propagates down the triangle: its first
entry is $1$ by the Gilbreath property, and every later entry is the
difference of two even entries.  Reading the last entry of each positive row
gives the claimed parity of the anti-diagonal; the bottom entry is $1$.
\end{proof}

\begin{lemma}[Zero-tail rigidity]\label{lem:zero-tail}
Let $S\in\G_n$ and $i\le n-2$.  If
\[
 e_{i+1}=e_{i+2}=\cdots=e_{n-2}=0,
\]
then $e_i\in\{0,2\}$.
\end{lemma}

\begin{proof}
Write row $i$ as $(x_1,\ldots,x_L)$.  Its first and last entries are
$x_1=1$ and $x_L=e_i$.  The displayed zero anti-diagonal entries are the
successive rightmost absolute differences of this row.  They force the
suffix $x_2,\ldots,x_L$ to be constant.  Indeed, once the last $m$ entries
are equal to $c$, the rightmost $m$-fold difference of
$(x_{L-m},c,\ldots,c)$ is $|x_{L-m}-c|$, so the next zero extends the
constant suffix.  (When $L=2$, the suffix assertion is immediate.)
Hence $x_2=e_i$.  The first entry of row $i+1$ is also $1$, and therefore
\[
 1=|x_2-x_1|=|e_i-1|.
\]
Thus $e_i$ is $0$ or $2$.
\end{proof}

\section{The sharp defective-width gap}

\begin{theorem}\label{thm:gap}
For every $n\ge3$ and every $S\in\G_n$,
\[
 K_S\ne C_S\quad\Longrightarrow\quad |K_S|\ge10.
\]
The constant $10$ is attained by $S=(2,3,5,9,15)$.
\end{theorem}

\begin{proof}
Suppose $K_S\ne C_S$.  By \eqref{eq:criterion}, an ordered completeness
inequality fails.  Let $i$ be the largest failing index and set
\[
 q=\sum_{j=i+1}^{n-1}e_j.
\]
All inequalities strictly to the right of $i$ hold.  Lemma~\ref{lem:parity}
shows that $q$ is odd.  If $q=1$, then
$e_{i+1}=\cdots=e_{n-2}=0$, so Lemma~\ref{lem:zero-tail} gives
$e_i\le2=1+q$, contrary to failure at $i$.  Hence
\begin{equation}\label{eq:qge3}
 q\ge3.
\end{equation}

Run the reverse recurrence through the tail to the right of $i$.  The
holding inequalities give
\begin{equation}\label{eq:tail}
 T_i=\{0,2,4,\ldots,q+1\}.
\end{equation}
For completeness, the induction starts with
$P_1(\{1\})=\{0,2\}$.  If the current set is
$\{0,2,\ldots,R\}$ and the next even entry $e$ satisfies $e\le R$, then
the permitted lower branch is $\{0,2,\ldots,e\}$ and the translated branch
is $\{e,e+2,\ldots,e+R\}$; their union is
$\{0,2,\ldots,e+R\}$.

At the failing index, $e_i>q+1$.  Thus every lower branch in
$P_{e_i}(T_i)$ is permitted.  The reflected lower copy and translated upper
copy meet only at $e_i$, so \eqref{eq:tail} yields
\begin{equation}\label{eq:five}
 |T_{i-1}|=2|T_i|-1=q+2\ge5,
 \qquad 0\notin T_{i-1}.
\end{equation}
Every remaining outer preimage step is cardinality-nondecreasing because it
contains the injective translate $e+T$.

If zero never enters, then $0\notin D_S$ and
\eqref{eq:reverse} gives $|K_S|=2|D_S|\ge10$.  Otherwise consider the first
outer step at which zero enters.  Its input $T$ has at least five elements
and does not contain zero.  The output can contain zero only if $e\in T$.
In that case $e+T$ consists of $|T|$ distinct positive elements, in addition
to zero.  The output has at least six elements, and subsequent steps cannot
decrease cardinality.  Hence
$|K_S|=2|D_S|-1\ge11$.  This proves the lower bound.

For sharpness, $(2,3,5,9,15)$ is Gilbreath and has ordered anti-diagonal
$(6,2,0,1)$.  Its reverse sets are
\[
 \{1\}\longmapsto\{0,2\}\longmapsto\{0,2\}
 \longmapsto\{0,2,4\}\longmapsto\{2,4,6,8,10\}.
\]
Reflection about $15$ gives
\[
 K_S=\{5,7,9,11,13,17,19,21,23,25\},
\]
which has size ten.  The candidate set also contains $15$, so it is indeed
defective.
\end{proof}

\section{Reduction of the width-at-most-nine problem}

\begin{corollary}\label{cor:reduction}
For $n\ge3$ and $S\in\G_n$,
\[
 |K_S|\le9
 \quad\Longleftrightarrow\quad
 K_S=C_S\ \text{ and }\ A(S)\le7.
\]
Equivalently, apart from the final entry $e_{n-1}=1$, the ordered
anti-diagonal has exactly one of the following forms; every unshown entry is
zero:
\begin{enumerate}
 \item $e_1=2$;
 \item $e_1=2$ and $e_q=2$ for one $2\le q\le n-2$;
 \item $e_1=2$ and $e_q=e_r=2$ for
       $2\le q<r\le n-2$;
 \item $e_1=4$ and $e_q=2$ for one $2\le q\le n-2$.
\end{enumerate}
The listed types have widths $5,7,9,9$, respectively, whenever they are
realized by a member of $\G_n$.
\end{corollary}

\begin{proof}
Theorem~\ref{thm:gap} rules out a defective set of width at most nine.  When
$K_S=C_S$, Lemma~\ref{lem:parity} makes $A(S)$ odd, and direct counting in
the parity interval gives $|C_S|=A(S)+2$.  This proves the first
equivalence.

Write $e_i=2a_i$ for $i<n-1$.  Strict increase of the top row gives
$a_1\ge1$, and
\[
 A(S)=1+2\sum_{i=1}^{n-2}a_i\le7,
 \qquad
 a_i\le1+\sum_{j=i+1}^{n-2}a_j
\]
where the second display is \eqref{eq:criterion} after division by two.
Thus $\sum a_i\le3$.  If the total is one, only type~1 occurs.  If it is
two, $a_1=2$ fails the first inequality, leaving type~2.  If the total is
three, the possibilities satisfying all ordered inequalities are either
$a_1=1$ with two later entries equal to one, or $a_1=2$ with one later entry
equal to one.  These are types~3 and~4.  Conversely, every listed type
satisfies the ordered inequalities, and its value of $A(S)$ gives the stated
width.
\end{proof}

\begin{remark}[What remains unresolved]
Corollary~\ref{cor:reduction} classifies only the possible ordered
anti-diagonals.  It neither decides which displayed tuples are realizable at
each length nor classifies all normalized top rows producing a realizable
tuple.  Those are precisely the missing steps needed for the original
if-and-only-if top-row classification.
\end{remark}

\section{Exact finite corroboration}

The proof above is independent of computation.  As a check, the ancillary
program \path{enumerate_threshold.py} exhaustively generated $\G_n$ for
$2\le n\le10$ in two ways: by the ordered reverse tree and by direct
difference-triangle tests.  For a positive append distance $d>A(S)$, all
folds occur without a sign change and the final value is $d-A(S)$; hence the
direct generator's scan of $1\le d\le A(S)+1$ is exhaustive.  The two exact
frontiers agreed at every length.  The resulting sequence counts were
\[
 1,1,2,6,27,180,1786,26094,559127
 \qquad(n=2,\ldots,10).
\]
For each $5\le n\le10$, the least defective width was ten, attained by
$1,1,8,9,37,102$ sequences, respectively; no defective sequence of width at
most nine occurred.  Every minimum-width hit was checked again by directly
constructing all appended triangles in the complete radius-$A(S)+1$ range.
All arithmetic used exact Python integers and finite sets.

The finite universe and the trust boundary of such a check should be stated
explicitly; a certificate-based treatment of this reporting principle in a
different finite combinatorial setting is given by Zeng, Liu, Ratnavelu, and
Huat \cite[Section~6]{ZengEtAl2026Certificates}.  Here the census is only
corroborative: no finite range is used to infer Theorem~\ref{thm:gap}.
It can be reproduced from the project root with
\begin{verbatim}
python3 evidence/enumerate_threshold.py --max-n 10 \
  --output /tmp/threshold_census_n10.json
\end{verbatim}

\section*{Computational and AI assistance disclosure}

OpenAI Codex assisted with literature searching, exact-computation design,
proof auditing, and preparation of the manuscript.  A separate model context
reconstructed the proof and independently recomputed the finite check through
$n=8$.  Such model passes are not peer review and are not mathematical
evidence.  The quantified result rests on the human-readable proof above;
the bounded computational claims rest on the exact program and recorded
output.  No proof assistant and no floating-point calculation were used.

\bibliographystyle{alpha}
\bibliography{references}

\end{document}
